\documentclass[12pt]{article}

\usepackage[dvips,letterpaper,margin=0.75in,bottom=0.75in]{geometry}
\usepackage{cite}
\usepackage{slashed}
\usepackage{graphicx}
\usepackage{mathtools}
\usepackage{amsmath}
\usepackage{amssymb}
\usepackage{braket}
%\usepackage[american,fulldiode]{circuitikz}
\usepackage{tikz}
\usepackage{arydshln}

\begin{document}

\title{PHY 130B \\ Lecture Notes E: \\
Electricity and Magnetism\\}
\author{Michael Mulhearn}

\maketitle

\newcounter{exercise}
\setcounter{exercise}{0}

\newcommand{\exercise}{
\stepcounter{exercise}
\noindent
{\bf Exercise E.\theexercise:}
}

\section{The Maxwell Equations}

Setting:
$$\hbar = c = \epsilon_0 = \mu_0 = 1$$
the Maxwell Equations become:
\begin{equation}
\label{eqn:maxwell}
\begin{aligned}
\nabla \cdot \vec{E} &= \rho \qquad \text{(a)} \qquad &\nabla \times \vec{E} &= -\frac{\partial \vec{B}}{\partial t} \qquad \text{(b)}\\[6pt]
\nabla \cdot \vec{B} &= 0 \qquad \text{(c)} \qquad &\nabla \times \vec{B} &= \vec{J} + \frac{\partial \vec{E}}{\partial t} \qquad \text{(d)}
\end{aligned}
\end{equation}
Taking the gradiant of (d):
\begin{equation*}
\nabla \cdot (\nabla \times \vec{B}) = \nabla \cdot \vec{J} + \frac{\partial}{\partial t} \, (\nabla \cdot \vec{E}) 
\end{equation*}
and then using (a) we have:
\begin{equation}
  \label{eqn:qcons-nabla}
\frac{\partial \rho}{\partial t} =  -\nabla \cdot \vec{J} 
\end{equation}
and so:
\begin{align*}
  \int_V \frac{\partial \rho}{\partial t} \; dV &=  \int_V -\nabla \cdot \vec{J} \; dV \\[7pt]
  \frac{d}{dt} \int_V \; \rho \; dV &=  \int_V \vec{J} \cdot dV 
\end{align*}
Assuming that $\vec{J}$ vanishes at infinity (or that no current leaves a finite volume), we have:
\begin{equation}
  \frac{d}{dt} \int_V \; \rho \; dV =  0 
\end{equation}
We see that electric charge conservation is a direct consequence of the Maxwell Equations.
Defining:
\begin{equation}
\label{eqn:four-current}  
[J^\mu] = (\rho, \vec{J})
\end{equation}
and recalling that:
\begin{equation*}
[\partial_\mu] = (\frac{\partial}{\partial t}, \nabla)
\end{equation*}
we can write Equation~\ref{eqn:qcons-nabla} as:
\begin{equation}
\label{eqn:conserved-four-current}
\partial_\mu J^\mu = 0
\end{equation}
and see that $J$ a {\it conserved four-current density}, but we most
often just call it a {\it conserved current}.  Since $\partial_\mu$ is
a four-vector and $0$ is a scalar, we see that $J^\mu$ is indeed a
four-vector.\\[5pt]

\exercise Consider a region of uniform charge density and no current so that in frame $F$:
$$[J^\mu]_F = (\rho, 0, 0, 0)$$\\
Using the Lorentz transformation, show that in frame $F'$ moving at
velocity $\beta$ in the $z$ direction,
$$[J^\mu]_F = (\gamma \rho, 0, 0, -\gamma \beta \rho)$$\\
Then explain these results in terms of length contraction (which
increases the charge density) and the appearance of current, both due
to the relative motion of $F'$ and $F$.

\section{Gauge Invariance of the Scalar and Vector Potentials}

As you will show in the exercises below, a field is curless if and
only if it arises from a conservative potential:
$$\nabla \times \vec{F} = 0 \iff \text{There exists} \; \phi \; \text{such that} \; \vec{F} = -\nabla \phi$$
In a similar manner, a field is curless if and only if it arises from the curl of a vector potential:
$$\nabla \cdot \vec{F} = 0 \iff \text{There exists} \; \vec{A} \; \text{such that} \; \vec{F} = \nabla \times \vec{A}$$
The existence of $\vec{A}$ is indeed gauranteed but quite challenging to
prove, so we will not attempt it here.  You will proof the other
direction in the exercises, however.

From Equation~\ref{eqn:maxwell}c of the Maxwell equations
$$\nabla \cdot \vec{B} = 0$$
we infer the existence of a vector potential $\vec{A}$ such that:
\begin{equation}
\label{eqn:bfield}
\vec{B} = \nabla \times \vec{A}
\end{equation}
The electric field is neither curlness nor divergenceless, however, starting from Equation~\ref{eqn:maxwell}b
$$\nabla \times \vec{E} \; = \; -\frac{\partial \vec{B}}{\partial t} \; = \; -\nabla \times \frac{\partial \vec{\vec{A}}}{\partial t} $$
so there is curless quantity:
$$\nabla \times \left( \vec{E} + \frac{\partial \vec{A}}{\partial t} \right) \; = \; 0$$
from which we infer the existence of a scalar potential:
$$\vec{E} + \frac{\partial \vec{A}}{\partial t} \; = \; -\nabla \phi$$
and so the electric field can be calculated from the scalar and vector potentials as:
\begin{equation}
\label{eqn:efield}
\vec{E} = -\nabla \phi - \frac{\partial \vec{A}}{\partial t}
\end{equation}

\noindent
Taken together, the scalar potential $\phi$ and the vector potenial
$\vec{A}$ uniquely specify the physical electromagentic fields
$\vec{E}$ and $\vec{B}$.  However, the potentials themselves are {\bf
  not} unique.  In what is the most important results of this section,
you will show that the electric and magnetic fields are unchanged by the transformation
$$\vec{A} \to \vec{A} + \nabla \lambda(\vec{x},t) \quad \text{and} \quad \phi \to \phi + \frac{\partial \lambda}{\partial t}$$
where $\lambda(\vec{x},t)$ is a real valued function of position and space.  We will show soon that the quanity:
$$[A_\mu] = (\phi, \vec{A})$$
is indeed a four vector, and so we can write this transformation as:
\begin{equation}
\label{eqn:gauge-transformation}
A_\mu \to A_\mu + \partial_\mu \, \lambda(x)
\end{equation}
where $x$ is the position-space four-vector.  We call this transformation a {\it gauge transformation}.\\[5pt]

\exercise Show that
\begin{equation*}
\nabla \cdot \left( \nabla \times \vec{F} \right) = 0 \qquad \text{and} \qquad \nabla \times \left(\nabla \phi \right) = 0
\end{equation*}
Use tensor notation with
$$(\nabla \times \vec{F})_i = \epsilon_{ijk}\partial_j F_k$$
and
$$\epsilon_{ijk} = \epsilon_{jki} = -\epsilon{ikj}. \qquad \square$$

\exercise Show that if
$$\nabla \times \vec{F} = 0$$
then
$$\phi(x) = \int_{x_0}^x \vec{F} \cdot d\vec{x}$$
is idependent of the path followed.  Hint: use Stoke's theorem.$\qquad \square$\\[5pt]

\exercise (*) Using Equation $\ref{eqn:bfield}$ show that the magnetic
field $\vec{B}$ is unchanged by the gauge transformation of Equation
$\ref{eqn:gauge-transformation}$.$\qquad \square$\\[5pt]

\exercise (*) Using Equation $\ref{eqn:efield}$ show that the electric
field $\vec{E}$ is unchanged by the gauge transformation of Equation
$\ref{eqn:gauge-transformation}$.$\qquad \square$\\[5pt]

\section{The Maxwell Equation in Lorentz Covariant Form}

Equations~\ref{eqn:bfield} and \ref{eqn:efield} express the
relationship between the electric and magnetic fields ($\vec{E}$ and
$\vec{B}$) and the scalar and vector potentials ($\phi$ and
$\vec{A}$).  You will verify below that these imply the homogenous
Maxwell Equations (Equations ~\ref{eqn:maxwell}a$\&$c) are already
satisfied.  In this section, we will express the remaining Maxwell
equations in terms of the four-potential $A^\mu$.

Starting from Equation~\ref{eqn:maxwell}a
$$\nabla \cdot \vec{E} = \rho$$
and inserting Equation~\ref{eqn:efield}
\begin{align*}
\nabla \cdot \left( -\nabla \phi - \frac{\partial A}{\partial t} \right) &= \rho \\
-\nabla^2\phi - \frac{\partial}{\partial t}(\nabla \cdot \vec{A}) &= \rho \\
\end{align*}
Our final results will look best when expressed in terms of the d'Alembertian operator
$$\square^2 = \partial^\mu \, \partial_\mu = \frac{\partial^2}{\partial t^2} - \nabla^2$$
which we shall call box squared.  Noting that:
$$-\nabla^2 = \square^2 - \frac{\partial^2}{\partial t^2}$$
we have:
\begin{equation}
\label{eqn:box-squared-phi}
  \Box^2 \phi - \frac{\partial}{\partial t}\left(\nabla \cdot \vec{A} + \frac{\partial \phi}{\partial t}\right) = \rho
\end{equation}
In a similar manner, starting from Equation~\ref{eqn:maxwell}d
\begin{equation*}
\nabla \times \vec{B} = \vec{J} + \frac{\partial \vec{E}}{\partial t}
\end{equation*}
you will show that:
\begin{equation}
\label{eqn:box-squared-A}
\Box^2 \vec{A} + \nabla\left(\nabla \cdot \vec{A} + \frac{\partial \phi}{\partial t}\right) = \vec{J}
\end{equation}
Paying close attention to the signs, note that:
$$\nabla \cdot \vec{A} + \frac{\partial \phi}{\partial t} =
\partial_\mu A^\mu$$ And so Equations~\ref{eqn:box-squared-phi} and
~\ref{eqn:box-squared-A} can be combined as:
\begin{equation}
\label{eqn:maxwell-covariant}
\Box^2 A^\mu + \partial^\mu( \partial_\nu A^\nu) = J^\mu  
\end{equation}
Together with Equations~\ref{eqn:efield} and \ref{eqn:bfield} this is
fully consistent with the original Maxwell equations.  This form is
manifestly Lorentz covariant, which shows that $A^\mu$ is
indeed a four-vector.  Often in this class, we have seen important
equations often take two forms: one for calculations
(e.g. Equations~\ref{eqn:maxwell}) and one for showing Lorentz
covariance explicitly (e.g. Equation~\ref{eqn:maxwell-covariant}.

Equation~\ref{eqn:maxwel-covariant} can be written equivalently as:
\begin{equation*}
\partial_\mu F^{\mu\nu} = J^\nu
\end{equation*}
where
\begin{equation}
F^{\mu\nu} = \partial^\mu A^\nu - \partial^\nu A^\mu
\end{equation}
is called the field strength tensor.  In the exercises, you will show that:
\begin{equation}
\label{eqn:field-strength-matrix}
  [F] = \begin{pmatrix}
    0 & -E_x & -E_y & -E_z \\
  E_x &  0   & -B_z & \phantom{-}B_y \\
  E_y &  \phantom{-}B_z &  0   & -B_x \\
  E_z & -B_y &  \phantom{-}B_x &  0   \\
  \end{pmatrix}
\end{equation}



The phenomenon of electricity and magnetism is a truly remarkable
bridge between classical physics and modern physics.  Electrical and
magnetic forces are part of our everyday lifes, and you can throughly
explore the phenomomena using magnets, circuits, and other things you
can hold in your own hands.  And yet, with no extra input, they are
consistent with special relativity.  It turns out, this is just the
beginning.  As you shall see, the E\&M bridge goes far indeed.\\[5pt]

\exercise Using Equation~\ref{eqn:bfield} verify that
Equation~\ref{eqn:maxwell}c is satisfied for any $A^\mu$.\\[5pt]

\exercise Using Equations~\ref{eqn:efield} and \ref{eqn:bfield},
verify that verify that Equation~\ref{eqn:maxwell}b is satisfied for
any $A^\mu$.\\[5pt]

\exercise Starting from Equation~\ref{eqn:maxwell}d and inserting
Equations~\ref{eqn:efield} and ~\ref{eqn:bfield}, derive Equation~\ref{eqn:box-squared-A}.\\[5pt]

\exercise (*) Verify the matrix representation in
Equation~\ref{eqn:field-strength-matrix}.  Here is a start:
$$F^{0i} = \frac{\partial A^i}{\partial t} + (\nabla \phi)^i = -E^i$$

\exercise Since you know how to transform tensors, you know how to
transform Electric and Magnetic fields now via
Equation~\ref{eqn:field-strength-matrix}.  For a Lorentz boost $L$ we
have:
$$[F]_F' = L [F]_F L$$
Write the field strength tensor $F$ for a uniform electric field of
strength $E$ in the $+z$ direction.  Apply a Lorentz boost of speed
$beta$ in the $+x$ direction.  Calculate the resulting electric and
magnetic fields.

\exercise Assume that the uniform electric field from the previous
problem is the result of a uniform sheet of charge with charge density
$\sigma$.  Show that result results from the previous calculation are
consistent with your expectations accounting for length contraction,
and the current introduced by the moving sheet.

\section{The Lorenz Gauge}

In the previous section, we saw that Maxwell's equations can be
written in covariant form as:
$$\Box^2 A^\mu + \partial^\mu( \partial_\nu A^\nu) = J^\mu$$  
and before that, we saw that electric and magetic fields are unchanged by a gauge transformation:
$$A_\mu \to A_\mu + \partial_\mu \, \lambda(x)$$
Because the physical fields are unchanged by a gauge transformation,
we are free to choose the gauge that we find most convenient.  The
Lorenz Gauge (different person than Lorentz!) is the condition:
$$\partial_\mu A^\mu = 0$$
We can always find a gauge transformation $\lambda$ that imposes this condition, because
$$\Box^2 \lambda(x) = -\partial_u A^\mu = f(x)$$
is simply a wave equation for a known source, which certainly has a
solution.

With the Lorenz gauge, the covariant form of the Maxwell equations takes its most elegant form:
\begin{equation}
\label{eqn:maxwell-lorenz}
\Box^2 A^\mu = J^\mu  
\end{equation}

\section{Global U(1) Symmetry}

By definition, the $U(1)$ symmetry group is the set of unitary $1 \times 1$ complex
matrices.  These are just the complex numbers of magnitude
one:
$$U(1) = \{ e^{i\alpha} \mid \alpha \in \mathcal{R}\}$$
because:
$$\left(e^{i\alpha}\right)^\dagger\, e^{i\alpha} = e^{-i\alpha} \, e^{i\alpha} = 1$$
Under a global $U(1)$ symmetry, the equations of motion are unchanged
by:
$$\psi \to e^{i\alpha} \psi$$
We already have several examples.  The Klein-Gordon equation
\begin{equation}
\label{eqn:klein-gordan}
(\Box^2 + m^2)\psi = 0
\end{equation}
is invariant under $U(1)$ symmetry because an overall phase simply cancels:
\begin{equation*}
  (\Box^2 + m^2) \left(e^{i\alpha} \psi \right) = 0
  \implies
  (\Box^2 + m^2) \psi = 0
\end{equation*}
The Schr\"odinger equation
\begin{equation}
\label{eqn:schrodinger}
i\frac{\partial\psi}{\partial t} = -\frac{\nabla^2}{2m}\psi
\end{equation}
and the Dirac equation
\begin{equation}
\label{eqn:dirac}
(i\gamma^\mu\partial_\mu - m)\psi = 0
\end{equation}
are each invariant under the $U(1)$ symmetry.  This implies that each
element of $U(1)$ commutes with Hamiltonian operator, and so the
generators of $U(1)$ will commute as well.

The generator of $U(1)$ is:
$$T = (-i) \frac{\partial}{\partial \alpha} e^{i\alpha} = 1$$
so we have:
$$\frac{d}{dt}\braket{1} = \frac{d}{dt}\braket{\psi \mid 1 \mid \psi} =
\frac{d}{dt} \braket{\psi \mid \psi} = 0$$
So the conserved quantity associated with a global $U(1)$ symmetry is
probability.  And, evidently, one is the generator of probability.

In the exercises, you will show that for the Dirac equation, the current:
\begin{equation}
\label{eqn:jmu-dirac}
j^\mu = \bar{\psi} \gamma^\mu \psi
\end{equation}
is conserved, so that:
\begin{equation}
\label{eqn:conserved-probability}
\partial_\mu j^\mu = 0.
\end{equation}
The time-like component is :
$$j^0 = \bar{\psi} \gamma^0 \psi = \psi^\dagger (\gamma^0)^2 \psi = \psi^\dagger \psi$$
and so:
$$\int \psi^\dagger \psi \, d^3x = \text{const}$$


\exercise Show that if
$$ \delta_\mu j^\mu = 0$$
in some (possibly infinite) volume $V$ and
$\vec{j}$ vanishes at the boundary, then
$$\frac{d}{dt} \int_V j^0 dV = 0$$

\exercise Show that Dirac Equation, can be written as:
$$\gamma^\mu \, \partial_\mu \psi = - i m \psi$$
and in the adjoint form:
$$\left( \partial_\mu \bar{\psi} \right) \, \gamma^\mu = + i m \bar{\psi}$$
Hint: you will need:
$$(\gamma^0)^2=1, \quad \gamma^0 \gamma^\mu \gamma^0 = \left( \gamma^\mu \right)^\dagger$$
~\\[-3pt]

\exercise Using the results of the previous problem, show that for Dirac equation:
$$\partial_\mu (\bar{\psi} \gamma^\mu \psi) = 0$$
~\\[-3pt]

\exercise Show that for the Klein-Gordon equation, the current:
\begin{equation*}
j^\mu = \frac{1}{2mi}\left(\psi^*\partial^\mu\psi - \psi\partial^\mu\psi^*\right)
\end{equation*}
satisfies $\partial_\mu j^\mu = 0$ using
Equation~\ref{eqn:klein-gordan}. Note however that $j^0$ is not
positive definite, and so cannot be interpreted as a probability
density.\\[5pt]

\exercise Show that for the Schr\"odinger equation, the probability density and current:
\begin{equation*}
\rho = \psi^*\psi \qquad \vec{j} = \frac{1}{2mi}\left(\psi^*\nabla\psi - \psi\nabla\psi^*\right)
\end{equation*}
satisfy the continuity equation:
\begin{equation*}
\frac{\partial\rho}{\partial t} + \nabla\cdot\vec{j} = 0
\end{equation*}
using Equation~\ref{eqn:schrodinger}.
~\\[-3pt]

\section{Charge Conservation}

For a Fermion of charge $Qe$, we can also consider the $U(1)$ symmetry
transformation:
$$\psi \to e^{i Qe \alpha} \psi$$
Our results only differ from the previous section by a factor of $Qe$,
and instead of a probability density we will have a charge density.
Now the generators is:
$$T = (-i) \frac{\partial}{\partial \alpha} e^{iQe\alpha} = Qe$$
with associated conserved current:
\begin{equation}
J^\mu = Q e j^\mu = Q e \bar{\psi} \gamma^\mu \psi
\end{equation}  
with
$$\partial_\mu J^\mu = 0$$
and:
$$J^0 = Q e \psi^\dagger \psi = \rho$$
is the charge density, which you should compare to
Equations~\ref{eqn:four-current} and \ref{eqn:conserved-four-current}.

\section{Local U(1) Symmetry}

In the previous section we saw that global $U(1)$ symmetry,
i.e. invariance under a global change of phase, is responsible for
probability conservation and (equivalently) charge conservation.
Next, let's consider what happens when we insist that our equations of
motion are unchanged by a {\it local} $U(1)$ symmetry.  That is, we
consider the transformation:
$$\psi \to e^{i Qe \lambda(x)} \psi$$
where $\lambda(x)$ is an arbitrary real function defined at each
space-time coordinate $x$.  Under this transformation, the Dirac equation:
\begin{align*}
  (i \gamma^\mu \partial_\mu - m ) \psi &= 0\\
  i \gamma^\mu \partial_\mu \psi &= m \psi 
\end{align*}
becomes:
\begin{align*}
  i \gamma^\mu \partial_\mu \left( e^{iQe\lambda} \psi \right) \; &= \; m \left( e^{iQe\lambda} \psi \right) \\[5pt]
  e^{iQe\lambda} \, i \gamma^\mu \left( \partial_\mu  + i Qe (\partial_\mu \lambda) \right)  \psi
  &= e^{iQe\lambda} \; m  \, \psi \\[5pt]
  i \gamma^\mu \left( \partial_\mu  + i Qe (\partial_\mu \lambda) \right)  \psi
  &= m  \, \psi 
\end{align*}
so the free-particle Dirac equation is not invariant under local $U(1)$ symmetry, as:
$$\partial_\mu \to \partial_\mu + i (\partial_\mu \lambda)$$
We can accommodate this symmetry by replacing $\partial_\mu$ with a new derivative:
$$D_\mu \equiv \partial_\mu + i \, Q \, e \, A_\mu$$
that adds a new term introducing a new interaction with a field $A_\mu(x)$. Now:
\begin{equation*}
D_\mu \left( e^{iQe\lambda(x)} \psi\right ) = e^{iQe\lambda(x)}
\left(\partial_\mu + i \, Q \, e \, (A_\mu + (\partial_\mu \lambda)) \right) \psi
\end{equation*}
If the new interaction is invariant under a gauge transformation:
$$A_\mu \to A_\mu + (\partial_\mu \lambda),$$
then we have:
\begin{equation*}
D_\mu \left( e^{iQe\lambda(x)} \psi \right) = e^{iQe\lambda(x)} D_\mu \, \psi
\end{equation*}
This shows that $D_\mu$ tranforms in the same manner as $\psi$, and so
we call $D_\mu$ a covariant derivative.
The Dirac equation using the covariant derivative $D_\mu$:
\begin{equation*}
  (i \gamma^\mu D_\mu - m ) \psi = 0
\end{equation*}
is invariant under the local $U(1)$ transformation:
The covariant derivative is not gauge invariant, and physical
quanitites must be guage invariant.  In the exercises, you will show
that the commutator of covariant derivatives is gauge invariant and:
\begin{equation}
\label{eqn:covariant-commutator}
[D_\mu, D_\nu] = i Q e (\partial_\mu A_\nu -\partial_\nu A_\mu).
\end{equation}
  
Thus, the symmetry requirements have compelled us to introduce a new
field $A_\mu$, and an interaction that is invariant under a gauge
transformation
$$A_\mu \to A_\mu + d_\mu \lambda$$
which is identical to that of the Maxwell equations.  
We can also construct a gauge invariant quantity for the interaction,
and it is identical to the electromagnetic field strength tensor:
$$[D_\mu, D_\nu] = i Q e F_{\mu\nu}$$
Using field theory, we would write down the most general
Lagrangian consistent with local $U(1)$ symmetry and derive the
equations of motion for the particles, obtaining the Dirac equation,
and the equations for $A^\mu$, which would be the Maxwell equations.  We
haven't shown things quite so directly here, but I hope you still find
these results compelling.  Local $U(1)$ symmetry requires an
interaction like electricity and magnetism.\\[5pt]

\exercise Verify Equation~\ref{eqn:covariant-commutator}.

\section{The Electromagnetic Interaction Hamiltonian}

We obtain the Hamiltonian from the Dirac Equation by noting that:
$$i \frac{d}{dt} \psi = H \psi$$
The Dirac equation, including the E+M interaction from $U(1)$ symmetry is:
$$i \gamma^\mu ( \partial_\mu + i Qe A_\mu) - m) \psi = 0$$









\end{document}
